\documentclass[11pt,a4paper]{article}
% Tema para Pre-Projeto de Mestrado (CIn/UFPE) - mesmo layout do formulario da coordenacao.
\usepackage[utf8]{inputenc}
\usepackage[T1]{fontenc}
\usepackage[brazil]{babel}
\usepackage{newtxtext}
\usepackage[scaled=0.95]{helvet}
\usepackage[a4paper,top=2.2cm,bottom=2.2cm,left=2.2cm,right=2.2cm]{geometry}
\usepackage{xcolor}
\usepackage{array}
\usepackage{ulem}
\normalem
\pagestyle{empty}
\setlength{\parindent}{1.25em}
\definecolor{faixa}{gray}{0.86}

\newcommand{\faixa}[2][\large]{%
  \par\noindent\colorbox{faixa}{\parbox{\dimexpr\linewidth-2\fboxsep}{\sffamily\bfseries#1 #2}}\par}

\begin{document}

\begin{center}
{\sffamily\large Universidade Federal de Pernambuco (UFPE)\\ Centro de Informática (CIn)}
\end{center}
\vspace{1.6em}

\noindent\colorbox{faixa}{\parbox{\dimexpr\linewidth-2\fboxsep}{\centering\rule{0pt}{1.6em}\sffamily\bfseries\Large Coordenação de Pós-Graduação em Ciência da Computação\rule[-0.7em]{0pt}{0pt}}}

\begin{center}
{\sffamily\bfseries\Large\uline{Tema para Pré-Projeto de Mestrado}}
\end{center}

\noindent
\begin{tabular}{@{\hspace{1em}}>{\raggedleft\bfseries}p{0.19\linewidth}@{\hspace{0.4em}}p{0.74\linewidth}}
Linha de pesquisa: & Visão Computacional, Realidade Virtual e Aumentada \\[2pt]
Título: & Integração de Dados Sintéticos e Reais para Recuperação Robusta de Cenas 4D \\[2pt]
Proponente: & Francisco Paulo Magalhães Simões \\
\end{tabular}

\vspace{0.8em}
\faixa{Descrição:}
\vspace{0.6em}

{\small
Métodos de recuperação de cenas dinâmicas (4D) baseados em Gaussian Splatting têm alcançado reconstruções fotorrealistas com renderização em tempo real [1]. Sua robustez, porém, ainda é limitada em condições frequentes em cenas reais, como movimentos rápidos, oclusões, poucas vistas e materiais com aparência complexa, e sua avaliação depende de bases reais pequenas, nas quais a geometria e o movimento de referência são difíceis de obter. Dados sintéticos gerados por motores de renderização oferecem controle total da cena e anotações exatas [2], mas a diferença de domínio em relação aos dados reais limita o seu uso direto [3]. Na replicação do Ref-Gaussian [4], realizada como etapa preliminar, o método atingiu cerca de 34 dB de PSNR em objetos sintéticos, enquanto o artigo original reporta entre 23 e 27 dB em cenas reais. Nesse contexto, esta pesquisa visa investigar o relacionamento entre dados sintéticos e reais para aumentar a robustez de métodos de recuperação de cenas 4D, analisando como gerar, combinar e ponderar esses dados no treinamento e na avaliação. O diferencial está na geração controlada de dados sintéticos dirigida às falhas observadas nos métodos atuais e na medição sistemática de quanto da robustez obtida se transfere para cenas reais.
\par}

\vspace{1.2em}
\faixa{Referências Bibliográficas:}
\vspace{0.6em}

{\small
[1] WU, Guanjun et al. 4D Gaussian Splatting for Real-Time Dynamic Scene Rendering. In: Proceedings of the IEEE/CVF Conference on Computer Vision and Pattern Recognition (CVPR), 2024.

[2] GREFF, Klaus et al. Kubric: A Scalable Dataset Generator. In: Proceedings of the IEEE/CVF Conference on Computer Vision and Pattern Recognition (CVPR), 2022.

[3] TOBIN, Josh et al. Domain Randomization for Transferring Deep Neural Networks from Simulation to the Real World. In: Proceedings of the IEEE/RSJ International Conference on Intelligent Robots and Systems (IROS), 2017.

[4] YAO, Yuxuan et al. Reflective Gaussian Splatting. In: Proceedings of the International Conference on Learning Representations (ICLR), 2025.
\par}

\end{document}
